% ============================================================
% SECCIÓN 3: Servicio de Borde — Raspberry Pi
% ============================================================
\chapter{Servicio de Borde: Raspberry Pi}
\label{sec:raspberry}

% ─────────────────────────────────────────────────────────────
\section{Misión y Posición en la Arquitectura del Sistema}

La Raspberry Pi ocupa el eslabón más crítico de toda la arquitectura Demeter: actúa como
\textbf{puerta de enlace entre el mundo físico y el mundo digital}. Sin ella, no existe ningún 
camino técnico para que la red de sensores y actuadores ESP32 ---que operan bajo un protocolo 
binario propietario de radiofrecuencia--- se comunique con el servidor central en la nube.

Su responsabilidad puede enunciarse de forma simple: \textit{escuchar la red de campo, 
comprender sus mensajes, y retransliterar esa información al lenguaje del servidor}; y 
a la inversa, \textit{recibir órdenes del servidor y traducirlas en señales tangibles para 
los nodos físicos}.

Esta posición intermedia, también llamada \textbf{Edge Gateway} o Capa de Borde, es 
especialmente relevante en arquitecturas IoT industriales porque evita que la lógica de 
protocolo de campo contamine el servidor, y que el servidor requiera acceso directo al hardware.

\begin{table}[H]
    \centering
    \caption{Tecnologías del servicio de borde}
    \begin{tabular}{lp{9cm}}
        \toprule
        \textbf{Componente} & \textbf{Tecnología y función} \\
        \midrule
        Lenguaje       & Python 3.9+ con modelo de concurrencia \texttt{asyncio} \\
        UART           & \texttt{pyserial-asyncio}: acceso no bloqueante al puerto serie \\
        Red            & WebSocket sobre TLS (\texttt{wss://}) via cable Ethernet \\
        Despliegue     & Contenedor Docker con mapeo de dispositivo físico al host \\
        Protocolo IoT  & Protocolo Demeter (propietario, binario, CRC-16) \\
        \bottomrule
    \end{tabular}
    \label{tab:rpi_techs}
\end{table}


% ─────────────────────────────────────────────────────────────
\section{Contenedorización y Acceso a Recursos del Sistema}

El servicio Python se ejecuta íntegramente dentro de un contenedor \textbf{Docker}, lo que garantiza 
un entorno reproducible y aislado del sistema operativo subyacente de la Raspberry Pi (Raspberry Pi OS).

El despliegue tiene tres características relevantes:

\begin{enumerate}
    \item \textbf{Mapeo del puerto serie}: El dispositivo físico \texttt{/dev/serial0} del host 
          es inyectado directamente dentro del contenedor mediante la directiva de mapeo de 
          dispositivos de Docker. Esto permite que el proceso Python en el contenedor lea y escriba 
          en el cable UART como si estuviera en la máquina física, sin pérdida de rendimiento ni 
          abstracción adicional.
    
    \item \textbf{Conectividad de red}: La Raspberry Pi utiliza exclusivamente \textbf{cable Ethernet} 
          para la conexión a internet. Esta decisión es deliberada: elimina la latencia y la 
          inestabilidad del Wi-Fi, especialmente relevante en entornos de invernadero con mucha 
          interferencia de RF. El contenedor hereda la red del host de forma directa.
    
    \item \textbf{Variables de entorno}: La URL del servidor (\texttt{WS\_URL}), la velocidad del 
          puerto serie (\texttt{BAUD\_RATE}), y los credenciales de la red son inyectados como 
          variables de entorno en tiempo de ejecución, permitiendo cambiar el entorno destino 
          (desarrollo, producción) sin recompilar la imagen.
\end{enumerate}

El Diagrama de Despliegue completo, mostrando la relación entre el host, el contenedor, el 
servidor y los nodos ESP32, se encuentra en la Sección~\ref{fig:anx_rpi_docker} del Anexo.


% ─────────────────────────────────────────────────────────────
\section{El Motor de Concurrencia: asyncio}

El servicio no emplea hilos del sistema operativo (\textit{threads}) para gestionar múltiples 
operaciones en paralelo. En su lugar, adopta el modelo de \textbf{concurrencia cooperativa} de 
Python mediante \texttt{asyncio}: un único hilo de ejecución gestiona múltiples tareas que 
voluntariamente ceden el control cuando están esperando una operación de entrada/salida.

Esto resulta especialmente eficiente para un servicio de este tipo, donde el 99\% del tiempo 
de espera es I/O (esperar bytes del puerto serie, o esperar mensajes de la red).

Las cuatro tareas concurrentes que el orquestador lanza al inicio son:

\begin{table}[H]
    \centering
    \caption{Tareas asyncio del servicio Raspberry Pi}
    \begin{tabular}{clp{7cm}}
        \toprule
        \# & \textbf{Tarea} & \textbf{Responsabilidad} \\
        \midrule
        1 & UART RX Loop      & Lee bytes del puerto serie de forma continua y no bloqueante. \\
        2 & UART TX Queue     & Consume de una cola FIFO interna y escribe tramas al puerto. \\
        3 & WebSocket Client  & Mantiene la conexión persistente con el servidor y envía telemetría. \\
        4 & WS Message Listener & Recibe y despacha mensajes entrantes (comandos) desde el servidor. \\
        \bottomrule
    \end{tabular}
    \label{tab:rpi_tasks}
\end{table}

El flujo completo del bucle de eventos y cómo se interconectan estas cuatro tareas a través de 
colas asíncronas puede consultarse en la Sección~\ref{fig:anx_rpi_asyncio} del Anexo.


% ─────────────────────────────────────────────────────────────
\section{El Módulo del Protocolo Demeter en la Raspberry Pi}

La Raspberry Pi no se limita a retransmitir bytes en bruto: \textbf{comprende el Protocolo Demeter}. 
Para ello, el servicio incorpora la misma lógica de serialización y deserialización descrita en el 
Capítulo~1, pero implementada en Python.

El módulo \textit{UART Processor} actúa como la frontera de descodificación. Cada byte recibido 
pasa por un \textit{pipeline} de validación:

\begin{enumerate}
    \item Detección del byte de sincronía de inicio (\texttt{0xAA}).
    \item Lectura de la cabecera: flags, nodo origen, nodo destino, tipo de comando y longitud.
    \item Verificación del \textit{Checksum} CRC-16 sobre el conjunto completo de la trama.
    \item Desglose del \textit{payload} según el tipo de comando reconocido.
    \item Disparo del evento interno correspondiente hacia las capas superiores.
\end{enumerate}

Si cualquiera de estos pasos falla, la trama es descartada silenciosamente y no se propaga ningún 
efecto, siguiendo la misma filosofía de fallo silente definida en el protocolo.

La información extraída de tramas válidas se clasifica según su tipo:

\begin{table}[H]
    \centering
    \caption{Tipos de trama reconocidos por el servicio Raspberry Pi}
    \begin{tabular}{llp{7.5cm}}
        \toprule
        \textbf{Tipo}      & \textbf{Nombre}   & \textbf{Acción desencadenada} \\
        \midrule
        \texttt{TH\_RPT}   & Sensor Report     & Publica datos de temperatura, humedad y suelo al servidor. \\
        \texttt{SYN}       & Handshake Step 1  & Responde con \texttt{SYN-ACK} y registra el nodo en el Device Manager. \\
        \texttt{ACK}       & Confirmación      & Marca la sesión como establecida para el nodo origen. \\
        \texttt{CMD\_ACK}  & Confirm. Actuador & Actualiza el estado del pin/actuador en la caché local. \\
        \texttt{NACK}      & Error de Nodo     & Notifica al servidor que el comando fue rechazado por el nodo. \\
        \bottomrule
    \end{tabular}
    \label{tab:rpi_frames}
\end{table}


% ─────────────────────────────────────────────────────────────
\section{Comunicación con el Backend: El Canal WebSocket}

La Raspberry Pi se conecta al servidor central usando un \textbf{WebSocket permanente y cifrado} 
sobre TLS (\texttt{wss://}), identificándose con el identificador reservado 
\texttt{raspberry\_gateway}. Este identificador es reconocido por el servidor como la fuente 
de verdad de la red de campo, y tiene privilegios especiales.

Esta decisión de diseño tiene implicaciones significativas:

\begin{designnote}
\textbf{Separación de canales}: El canal WebSocket es exclusivo para comunicaciones máquina a 
máquina (Raspberry $\leftrightarrow$ Backend). Los clientes del navegador (Frontend) usan el mismo protocolo 
WebSocket, pero mediante un identificador diferente y únicamente en modo \textit{recepción}. 
Los comandos desde el navegador se envían como peticiones HTTP normales; el servidor los 
convierte en mensajes WebSocket hacia la Raspberry internamente. Esto simplifica el frontend 
al eliminar la lógica de reconexión y autenticación de sockets en el navegador.
\end{designnote}

La conexión WebSocket es bidireccional:
\begin{itemize}[leftmargin=*, label=--]
    \item \textbf{Dirección ascendente (RPi → Servidor)}: La Raspberry publica eventos JSON 
          para cada trama de telemetría recibida, confirmaciones de actuadores y nuevos nodos.
    \item \textbf{Dirección descendente (Servidor → RPi)}: El servidor envía comandos JSON 
          (set\_gpio, exec\_sequence, ping) que el Command Dispatcher traduce a tramas binarias.
\end{itemize}

El flujo ascendente se ilustra en la Sección~\ref{fig:anx_rpi_telemetria} del Anexo.
El flujo descendente se ilustra en la Sección~\ref{fig:anx_rpi_comandos} del Anexo.


% ─────────────────────────────────────────────────────────────
\section{El Command Dispatcher: Adaptador JSON $\leftrightarrow$ Binario}

El componente más delicado de la Raspberry Pi para el flujo de comandos es el 
\textbf{Command Dispatcher}. Su rol es exactamente opuesto al del \textit{UART Processor}: 
mientras este último convierte tramas binarias en estructuras de datos de alto nivel, 
el Dispatcher hace el camino inverso.

Cuando el servidor envía un comando JSON como:
\begin{lstlisting}[style=trama, caption={Ejemplo de comando JSON recibido por la Raspberry Pi}]
{
  "type": "set_gpio",
  "target_id": 3,
  "pin": 4,
  "value": 1,
  "duration_ms": 15000
}
\end{lstlisting}

El Dispatcher realiza los siguientes pasos:

\begin{enumerate}
    \item \textbf{Validación de esquema}: Verifica que el JSON contenga todos los campos 
          requeridos para el tipo de comando indicado (\texttt{type}).
    \item \textbf{Resolución de destino}: Consulta al Device Manager para verificar 
          que el nodo \texttt{target\_id} está registrado y activo.
    \item \textbf{Empaquetación del payload}: Serializa los parámetros del comando 
          (\texttt{pin}, \texttt{value}, \texttt{duration\_ms}) siguiendo el formato 
          binario del Protocolo Demeter.
    \item \textbf{Construcción de la cabecera}: Prepende el byte de sincronía, 
          los identificadores de origen y destino, el código de comando y la longitud.
    \item \textbf{Cálculo del CRC-16}: Genera el checksum de todo el paquete.
    \item \textbf{Encola en la TX Queue}: Inserta los bytes finales en la cola 
          de transmisión del módulo UART, desde donde la tarea \textit{TX Loop} los 
          enviará al puerto serie en la siguiente iteración del bucle de eventos.
\end{enumerate}


% ─────────────────────────────────────────────────────────────
\section{Device Manager: La Caché Local de Topología}

Para no depender de una consulta al servidor en cada operación, la Raspberry Pi mantiene 
\textbf{en memoria una tabla de estado local} de todos los nodos conocidos en la red.

El \textit{Device Manager} es el módulo responsable de esta caché. Registra:
\begin{itemize}[leftmargin=*, label=--]
    \item \textbf{Identidad del nodo}: identificador numérico y dirección MAC.
    \item \textbf{Instante de conexión}: timestamp del último mensaje recibido (\textit{last\_seen}).
    \item \textbf{Nivel de batería}: último valor reportado de tensión de la batería.
    \item \textbf{Estado de pines}: último estado conocido (ON/OFF) de cada salida de actuación.
\end{itemize}

Esta información permite al Dispatcher verificar si un nodo está activo antes de enviar 
un comando, y al \textit{UART Processor} determinar si una trama SYN proviene de un nodo 
previamente conocido o de uno nuevo que requiere registro. La persistencia de esta tabla 
en disco (archivo \texttt{devices.json}) garantiza que tras un reinicio del contenedor, 
el servicio no pierde la memoria de los nodos con los que ya ha establecido sesión.
